\documentclass[11pt,a4paper]{article}
% preamble.tex — shared setup for human-verification documents
% Compile each document with:  tectonic <File>.tex

\usepackage[T1]{fontenc}
\usepackage[utf8]{inputenc}
\usepackage{lmodern}
\usepackage{amssymb,amsmath,amsthm}
\usepackage{mathtools}
\usepackage[dvipsnames]{xcolor}
\usepackage{listings}
\usepackage{hyperref}
\usepackage[margin=2.5cm]{geometry}
\usepackage{parskip}

% ── Colours ──────────────────────────────────────────────────────────
\definecolor{leanBg}{HTML}{F7F7F0}
\definecolor{haskellBg}{HTML}{F0F4F7}
\definecolor{keyword}{HTML}{7B3294}
\definecolor{comment}{HTML}{5A8C32}
\definecolor{noteFg}{HTML}{B35900}

% ── Lean 4 listings language ─────────────────────────────────────────
% NOTE: Lean listing blocks use ASCII approximations of unicode symbols
% (e.g. -> for arrow, o. for compose).  The literate key substitutes
% these with proper math-mode renderings in the PDF.
\lstdefinelanguage{Lean4}{
  morekeywords={def,theorem,structure,where,namespace,end,import,fun,
                let,have,show,by,exact,intro,cases,rfl,subst,obtain,
                funext,simp,rw,apply,variable,Type,Prop,instance,
                section,open,class,extends,deriving,sorry,noncomputable},
  sensitive=true,
  morecomment=[l]{--},
  morecomment=[n]{/-}{-/},
  morestring=[b]",
}

% ── Lean listing style ───────────────────────────────────────────────
\lstdefinestyle{lean}{
  language=Lean4,
  backgroundcolor=\color{leanBg},
  basicstyle=\small\ttfamily,
  keywordstyle=\color{keyword}\bfseries,
  commentstyle=\color{comment}\itshape,
  breaklines=true,
  frame=single,
  framerule=0.4pt,
  xleftmargin=4pt,
  xrightmargin=4pt,
  aboveskip=6pt,
  belowskip=4pt,
  literate=
    {->}{{$\to$}}2
    {=>}{{$\Rightarrow$}}2
    {<|}{{$\langle$}}1
    {|>}{{$\rangle$}}1
    {==>}{{$\Longrightarrow$}}3
    { o. }{{$\;\circ\;$}}3
    {forall}{{$\forall$}}1
    {:=}{{$\,:=\,$}}2
    {!=}{{$\neq$}}2
    {<=}{{$\leq$}}2
    {>=}{{$\geq$}}2,
}

% ── Haskell listing style ───────────────────────────────────────────
\lstdefinestyle{haskell}{
  language=Haskell,
  backgroundcolor=\color{haskellBg},
  basicstyle=\small\ttfamily,
  keywordstyle=\color{keyword}\bfseries,
  commentstyle=\color{comment}\itshape,
  breaklines=true,
  frame=single,
  framerule=0.4pt,
  xleftmargin=4pt,
  xrightmargin=4pt,
  aboveskip=6pt,
  belowskip=4pt,
  literate=
    {->}{{$\to$}}2
    {=>}{{$\Rightarrow$}}2,
}

% ── Theorem environments ────────────────────────────────────────────
\theoremstyle{definition}
\newtheorem{defn}{Definition}[section]
\newtheorem{thm}[defn]{Theorem}
\newtheorem{prop}[defn]{Proposition}

\theoremstyle{remark}
\newtheorem*{fidelitynote}{Fidelity note}

% ── Convenience macros ───────────────────────────────────────────────
\newcommand{\lean}[1]{\texttt{#1}}      % inline Lean identifier
\newcommand{\hask}[1]{\texttt{#1}}      % inline Haskell identifier
\newcommand{\Ext}{\mathrm{Ext}}
\newcommand{\Pos}{\mathrm{Pos}}
\newcommand{\Shape}{\mathrm{Shape}}


\title{Human Verification: \texttt{Trust}\\[4pt]
  \large The Certificate Lemma and Transport of Trust}
\author{Machine-generated from \texttt{lean/Trust.lean} for review by Neil Ghani}
\date{July 2026}

\begin{document}
\maketitle

\begin{abstract}
This document presents every definition and theorem \emph{statement} from the
Lean~4 module \lean{Containers.Trust} in three parallel forms: the original
Lean, a readable Haskell rendering, and precise mathematical English.  Proofs
are omitted --- they are machine-checked by Lean.  This is the central file
of the trust-systems formalisation, providing the definitions and results
referenced by all other modules.

Main results formalised here: (i) the certificate lemma (Lemma 2.2): $\tau$
is the greatest sound assignment; (ii) equivalence of the fixpoint form and
the derivation-sup (max-min) form of $\tau$; (iii) lax transport of trust
along morphisms (Lemma 3.2); (iv) consolidation preservation (Theorem 3):
a genuine position map preserves $\tau$ exactly; (v) the circular-citation
counterexample; and (vi) the shape-only counterexample.
\end{abstract}

\tableofcontents
\bigskip

%===================================================================
\section{Overview}
%===================================================================

A \emph{trust system} is a dependency structure graded in a poset of check
strengths $T$.  Claims have checking steps with premises; the grade records
how strong a check each step received.  The \emph{canonical trust} $\tau(x)$
is the widest-bottleneck value: the best derivation's weakest link.

Two corrections to the prose notes discovered during formalisation:

\begin{enumerate}
\item Lemma 2.2 (i) needs $T$ to be a \textbf{chain} (linear order), not
  merely a lattice: in a general lattice, a finite sup need not be attained
  by any element.  All deployed check-posets are finite chains.
\item Lemma 3.2 needs $\phi(\bot) = \bot$ (``unchecked maps to unchecked''):
  monotonicity alone fails on underivable claims.
\end{enumerate}

%===================================================================
\section{The check-strength chain}
%===================================================================

\subsection{Class: \lean{Chain}}

\begin{defn}[Chain]
A \textbf{chain} (bounded decidable linear order) is a type $T$ with a
relation $\le$, a bottom element $\bot$, satisfying: reflexivity,
transitivity, antisymmetry, totality, $\bot \le a$ for all $a$, and
decidable $\le$.
\end{defn}

\paragraph{Lean.}
\begin{lstlisting}[style=lean]
class Chain (T : Type) extends LE T where
  bot       : T
  le_refl   : forall a : T, a <= a
  le_trans  : forall {a b c : T}, a <= b -> b <= c -> a <= c
  le_antisymm : forall {a b : T}, a <= b -> b <= a -> a = b
  le_total  : forall a b : T, a <= b \/ b <= a
  bot_le    : forall a : T, bot <= a
  dec       : forall a b : T, Decidable (a <= b)
\end{lstlisting}

\paragraph{Haskell.}
\begin{lstlisting}[style=haskell]
class (Eq t, Ord t, Bounded t) => Chain t where
  bot :: t
  bot = minBound
-- In Haskell, Ord + Bounded gives a total order with a bottom.
-- Decidability is automatic (all Haskell values are decidable).
\end{lstlisting}

\begin{fidelitynote}
In Lean the chain structure is a typeclass with explicit proof fields.
Lean's \lean{instance} mechanism checks these proofs at class instantiation
time.  In Haskell, \hask{Ord} provides totality and decidability; antisymmetry
and transitivity are assumed (no proof terms).
\end{fidelitynote}

%===================================================================
\section{Meet and join on a chain}
%===================================================================

\subsection{Definitions: \lean{cmin} and \lean{cmax}}

\begin{defn}[Meet and join]
For $a, b : T$ in a chain:
\begin{align*}
  \mathrm{cmin}(a, b) &= \text{if } a \le b \text{ then } a \text{ else } b
    && (\text{meet, weaker of the two}), \\
  \mathrm{cmax}(a, b) &= \text{if } a \le b \text{ then } b \text{ else } a
    && (\text{join, stronger of the two}).
\end{align*}
\end{defn}

\paragraph{Lean.}
\begin{lstlisting}[style=lean]
def cmin (a b : T) : T := if a <= b then a else b
def cmax (a b : T) : T := if a <= b then b else a
\end{lstlisting}

\paragraph{Haskell.}
\begin{lstlisting}[style=haskell]
cmin :: (Ord t) => t -> t -> t
cmin a b = if a <= b then a else b

cmax :: (Ord t) => t -> t -> t
cmax a b = if a <= b then b else a

-- Equivalently: cmin = min, cmax = max.
\end{lstlisting}

\subsection{Theorems: \lean{cmin} and \lean{cmax} properties}

The following basic facts about \lean{cmin} and \lean{cmax} are proved
and used throughout:

\begin{thm}[Meet and join properties]
\begin{align*}
  &\mathrm{cmin}(a,b) \le a && \text{(\lean{cmin\_le\_left})} \\
  &\mathrm{cmin}(a,b) \le b && \text{(\lean{cmin\_le\_right})} \\
  &c \le a \wedge c \le b \implies c \le \mathrm{cmin}(a,b) && \text{(\lean{le\_cmin})} \\
  &a \le \mathrm{cmax}(a,b) && \text{(\lean{le\_cmax\_left})} \\
  &b \le \mathrm{cmax}(a,b) && \text{(\lean{le\_cmax\_right})} \\
  &\mathrm{cmax}(a,b) = a \vee \mathrm{cmax}(a,b) = b && \text{(\lean{cmax\_cases})}
\end{align*}
The last property (\lean{cmax\_cases}) uses linearity of the chain and is the
attainment fact that makes the soundness direction of the certificate lemma work.
\end{thm}

\paragraph{Lean.}
\begin{lstlisting}[style=lean]
theorem cmin_le_left (a b : T) : cmin a b <= a
theorem cmin_le_right (a b : T) : cmin a b <= b
theorem le_cmin {a b c : T} (ha : c <= a) (hb : c <= b) : c <= cmin a b
theorem le_cmax_left (a b : T) : a <= cmax a b
theorem le_cmax_right (a b : T) : b <= cmax a b
theorem cmax_cases (a b : T) : cmax a b = a \/ cmax a b = b
\end{lstlisting}

\paragraph{Haskell.}
\begin{lstlisting}[style=haskell]
-- Standard properties of min/max on a total order.
-- cmax_cases: max a b == a || max a b == b   (always True)
\end{lstlisting}

%===================================================================
\section{Folds over lists}
%===================================================================

\subsection{Theorems: folds of \lean{cmin} and \lean{cmax} over lists}

The following are helper lemmas used in the proofs of the main theorems.
They express standard properties of iterated meets and joins over lists.

\begin{thm}[List fold lemmas]
Let $l$ be a list and $\mathrm{init}$ an initial value.
\begin{align*}
  &a \in l \implies a \le l.\mathrm{foldr}(\mathrm{cmax}, \mathrm{init})
  && \text{(\lean{le\_foldr\_max})} \\
  &l.\mathrm{foldr}(\mathrm{cmax}, \mathrm{init}) = \mathrm{init}
   \;\vee\; l.\mathrm{foldr}(\mathrm{cmax}, \mathrm{init}) \in l
  && \text{(\lean{foldr\_max\_cases})} \\
  &a \in l \implies l.\mathrm{foldr}(\mathrm{cmin}, \mathrm{init}) \le a
  && \text{(\lean{foldr\_min\_le})} \\
  &l.\mathrm{foldr}(\mathrm{cmin}, \mathrm{init}) \le \mathrm{init}
  && \text{(\lean{foldr\_min\_le\_init})}
\end{align*}
\lean{foldr\_max\_cases} is the key attainment lemma: on a chain, the
max-fold is always attained by the initial value or a list member.
\end{thm}

\paragraph{Lean.}
\begin{lstlisting}[style=lean]
theorem le_foldr_max {a : T} {l : List T} (h : a in l) (init : T) :
    a <= l.foldr cmax init

theorem foldr_max_cases : forall (l : List T) (init : T),
    l.foldr cmax init = init \/ l.foldr cmax init in l

theorem foldr_min_le {a : T} {l : List T} (h : a in l) (init : T) :
    l.foldr cmin init <= a

theorem foldr_min_le_init : forall (l : List T) (init : T),
    l.foldr cmin init <= init
\end{lstlisting}

\paragraph{Haskell.}
\begin{lstlisting}[style=haskell]
-- le_foldr_max:    a `elem` l => a <= foldr cmax init l
-- foldr_max_cases: foldr cmax init l == init || foldr cmax init l `elem` l
-- foldr_min_le:    a `elem` l => foldr cmin init l <= a
-- foldr_min_le_init: foldr cmin init l <= init
\end{lstlisting}

\subsection{Theorems: monotonicity lemmas for folds}

\begin{thm}[Monotonicity of folds]
\begin{enumerate}
\item \textbf{\lean{foldr\_min\_map\_mono}}: pointwise domination of mapped
  \lean{cmin}-folds: if $f(a) \le g(a)$ for all $a \in l$ and $i_1 \le i_2$,
  then $(l.\mathrm{map}\,f).\mathrm{foldr}(\mathrm{cmin}, i_1) \le
  (l.\mathrm{map}\,g).\mathrm{foldr}(\mathrm{cmin}, i_2)$.

\item \textbf{\lean{phi\_foldr\_min}}: a monotone map is lax for
  \lean{cmin}-folds:
  $\phi(l.\mathrm{foldr}(\mathrm{cmin}, \mathrm{init})) \le
  (l.\mathrm{map}\,\phi).\mathrm{foldr}(\mathrm{cmin}, \phi(\mathrm{init}))$.
  Holds for any monotone $\phi$ without meet-preservation.

\item \textbf{\lean{mapCongr}}: congruence for \lean{List.map} on members:
  if $f(a) = g(a)$ for all $a \in l$ then $l.\mathrm{map}\,f = l.\mathrm{map}\,g$.
\end{enumerate}
\end{thm}

\paragraph{Lean.}
\begin{lstlisting}[style=lean]
theorem foldr_min_map_mono {alpha : Type} {f g : alpha -> T} {i1 i2 : T} :
    forall (l : List alpha),
    (forall a, a in l -> f a <= g a) -> i1 <= i2 ->
    (l.map f).foldr cmin i1 <= (l.map g).foldr cmin i2

theorem phi_foldr_min {T' : Type} [Chain T'] {phi : T -> T'}
    (hmono : forall {a b : T}, a <= b -> phi a <= phi b) :
    forall (l : List T) (init : T),
    phi (l.foldr cmin init) <= (l.map phi).foldr cmin (phi init)

theorem mapCongr {alpha beta : Type} {f g : alpha -> beta} :
    forall (l : List alpha),
    (forall a, a in l -> f a = g a) -> l.map f = l.map g
\end{lstlisting}

\paragraph{Haskell.}
\begin{lstlisting}[style=haskell]
-- foldr_min_map_mono: pointwise f<=g, i1<=i2 =>
--   foldr cmin i1 (map f l) <= foldr cmin i2 (map g l)
-- phi_foldr_min: phi monotone =>
--   phi (foldr cmin init l) <= foldr cmin (phi init) (map phi l)
-- mapCongr: (forall a in l, f a == g a) => map f l == map g l
\end{lstlisting}

%===================================================================
\section{Trust systems}
%===================================================================

\subsection{Structure: \lean{System}}

\begin{defn}[Trust system]
A \textbf{trust system} over claims $C$ graded in chain $T$ consists of:
\begin{itemize}
  \item a type \lean{Step} of checking steps,
  \item $\mathrm{concl} : \lean{Step} \to C$ (the conclusion of each step),
  \item $\mathrm{prems} : \lean{Step} \to \lean{List}\,C$ (finite premises),
  \item $\mathrm{grade} : \lean{Step} \to T$ (the check strength received),
  \item $\mathrm{stepsInto} : C \to \lean{List}\,\lean{Step}$ (all steps
    concluding that claim --- finite, as the registry is a finite artifact),
  \item coherence: $f \in \mathrm{stepsInto}(x) \iff \mathrm{concl}(f) = x$.
\end{itemize}
\end{defn}

\paragraph{Lean.}
\begin{lstlisting}[style=lean]
structure System (C : Type) (T : Type) [Chain T] where
  Step       : Type
  concl      : Step -> C
  prems      : Step -> List C
  grade      : Step -> T
  stepsInto  : C -> List Step
  mem_stepsInto : forall {x : C} {f : Step},
    f in stepsInto x -> concl f = x
  stepsInto_mem : forall f : Step, f in stepsInto (concl f)
\end{lstlisting}

\paragraph{Haskell.}
\begin{lstlisting}[style=haskell]
data System c t = System
  { step        :: [()]          -- just the count; we index by ()
  , concl       :: () -> c
  , prems       :: () -> [c]
  , grade       :: () -> t
  , stepsInto   :: c -> [()]
  -- coherence conditions are Prop in Lean, omitted in Haskell
  }
-- In practice: Step is a concrete type like Bool or Evt.
\end{lstlisting}

\begin{fidelitynote}
In Lean, \lean{Step} is a field of the structure and can vary between
instances.  The Haskell approximation glosses over this type-level variation;
in practice each instance uses a concrete \hask{Step} type (e.g., \hask{Bool}
for the cyclic example, \hask{Evt} for trajectories).
\end{fidelitynote}


\subsection{Definition: \lean{System.Prem}}

\begin{defn}[Premise relation]
$y$ is a \textbf{premise} of $x$ in system $S$ if there exists a checking
step $f \in \mathrm{stepsInto}(x)$ with $y \in \mathrm{prems}(f)$:
\[
  S.\mathrm{Prem}(y, x) \;\iff\;
  \exists f,\; f \in \mathrm{stepsInto}(x) \wedge y \in \mathrm{prems}(f).
\]
Well-foundedness of this relation is the hypothesis of the certificate lemma.
It is \emph{load-bearing}: the circular counterexample (\lean{cyc}) has
$\neg\,\mathrm{WellFounded}(\mathrm{Prem})$.
\end{defn}

\paragraph{Lean.}
\begin{lstlisting}[style=lean]
def System.Prem (y x : C) : Prop :=
  exists f, f in S.stepsInto x /\ y in S.prems f
\end{lstlisting}

\paragraph{Haskell.}
\begin{lstlisting}[style=haskell]
prem :: System c t -> c -> c -> Bool
prem s y x = any (\f -> y `elem` prems s f) (stepsInto s x)
\end{lstlisting}


\subsection{Definition: \lean{stepGrade}}

\begin{defn}[Step grade under an assignment]
Given an assignment $t : C \to T$, the \textbf{grade} of step $f$ under $t$
is the meet of the step's grade with the assigned levels of all its premises:
\[
  \mathrm{stepGrade}(t, f) =
  \bigl(\mathrm{prems}(f).\mathrm{map}(t)\bigr).\mathrm{foldr}(\mathrm{cmin},\;
  \mathrm{grade}(f)).
\]
This is the ``weakest link'' value: the step is only as strong as its weakest
premise's assigned trust, bounded by the step's own check grade.
\end{defn}

\paragraph{Lean.}
\begin{lstlisting}[style=lean]
def stepGrade (t : C -> T) (f : S.Step) : T :=
  ((S.prems f).map t).foldr cmin (S.grade f)
\end{lstlisting}

\paragraph{Haskell.}
\begin{lstlisting}[style=haskell]
stepGrade :: System c t -> (c -> t) -> StepOf s -> t
stepGrade s t f = foldr cmin (grade s f) (map t (prems s f))
\end{lstlisting}

%===================================================================
\section{Canonical trust}
%===================================================================

\subsection{Definition: \lean{tau}}

\begin{defn}[Canonical trust (fixpoint form)]
Given a well-founded premise relation, the \textbf{canonical trust}
$\tau : C \to T$ is defined by well-founded recursion:
\[
  \tau(x) = \bigvee_{f \in \mathrm{stepsInto}(x)} \mathrm{stepGrade}(\tau, f),
\]
where $\bigvee$ is the \lean{cmax}-fold with initial value $\bot$.
The internal definition uses a \lean{dite} guard defaulting to $\bot$ off the
premise relation (erased by \lean{tau\_eq}).
\end{defn}

\paragraph{Lean.}
\begin{lstlisting}[style=lean]
noncomputable def tau (wf : WellFounded S.Prem) : C -> T :=
  wf.fix fun x rec =>
    ((S.stepsInto x).map fun f =>
      ((S.prems f).map fun y =>
        if h : S.Prem y x then rec y h else Chain.bot).foldr
          cmin (S.grade f)).foldr
      cmax Chain.bot
\end{lstlisting}

\paragraph{Haskell.}
\begin{lstlisting}[style=haskell]
-- tau s x = foldr cmax bot
--             (map (\f -> foldr cmin (grade s f)
--                           (map (tau s) (prems s f)))
--               (stepsInto s x))
-- Non-computable in Lean (Classical); in Haskell we assume termination.
tau :: System c t -> c -> t
tau s x = foldr cmax bot
  [foldr cmin (grade s f) (map (tau s) (prems s f))
  | f <- stepsInto s x]
\end{lstlisting}


\subsection{Theorem: \lean{tau\_eq}}

\begin{thm}[Clean unfolding of $\tau$]
\[
  \tau(x) = \bigl(\mathrm{stepsInto}(x).\mathrm{map}(\mathrm{stepGrade}(\tau,
  {-}))\bigr).\mathrm{foldr}(\mathrm{cmax}, \bot).
\]
This is the intended fixpoint equation; the \lean{dite} guard in the
definition collapses to this on well-founded input.
\end{thm}

\paragraph{Lean.}
\begin{lstlisting}[style=lean]
theorem tau_eq (wf : WellFounded S.Prem) (x : C) :
    tau S wf x =
      ((S.stepsInto x).map (stepGrade S (tau S wf))).foldr
        cmax Chain.bot
\end{lstlisting}

\paragraph{Haskell.}
\begin{lstlisting}[style=haskell]
-- tau_eq: tau s x = foldr cmax bot (map (stepGrade s (tau s)) (stepsInto s x))
\end{lstlisting}

%===================================================================
\section{Soundness: the certificate lemma (Lemma 2.2)}
%===================================================================

\subsection{Definition: \lean{Sound}}

\begin{defn}[Sound assignment]
An assignment $t : C \to T$ is \textbf{sound} if every claim trusted above
$\bot$ has a node-local certificate: a checking step whose \lean{stepGrade}
under $t$ dominates $t(x)$:
\[
  \mathrm{Sound}(t) \;\iff\;
  \forall x : C,\; t(x) \ne \bot \;\Rightarrow\;
  \exists f \in \mathrm{stepsInto}(x),\; t(x) \le \mathrm{stepGrade}(t, f).
\]
This is precisely the deployed boundary rule: the validator checks exactly
this predicate at each node.
\end{defn}

\paragraph{Lean.}
\begin{lstlisting}[style=lean]
def Sound (t : C -> T) : Prop :=
  forall x : C, t x != Chain.bot ->
    exists f, f in S.stepsInto x /\ t x <= stepGrade S t f
\end{lstlisting}

\paragraph{Haskell.}
\begin{lstlisting}[style=haskell]
sound :: System c t -> (c -> t) -> Bool
sound s t = all check (claims s)
  where
    check x = t x == bot ||
              any (\f -> t x <= stepGrade s t f) (stepsInto s x)
\end{lstlisting}


\subsection{Theorem: \lean{sound\_le\_tau} (Lemma 2.2 ii)}

\begin{thm}[Every sound assignment is bounded by $\tau$]
If $t$ is sound, then $t(x) \le \tau(x)$ for all $x$:
\[
  \mathrm{Sound}(t) \;\Rightarrow\; \forall x,\; t(x) \le \tau(x).
\]
Node-local certificates imply the global bound; claiming above $\tau$ is
trust inflation and must fail the certificate somewhere.
\end{thm}

\paragraph{Lean.}
\begin{lstlisting}[style=lean]
theorem sound_le_tau (wf : WellFounded S.Prem) {t : C -> T}
    (hs : Sound S t) : forall x, t x <= tau S wf x
\end{lstlisting}

\paragraph{Haskell.}
\begin{lstlisting}[style=haskell]
-- sound_le_tau: sound s t => forall x, t x <= tau s x
\end{lstlisting}


\subsection{Theorem: \lean{tau\_sound} (Lemma 2.2 i)}

\begin{thm}[$\tau$ is sound]
The canonical trust is itself sound:
\[
  \mathrm{Sound}(\tau).
\]
This uses \lean{foldr\_max\_cases} (attainment), hence linearity of the
chain.  Together with \lean{sound\_le\_tau}: $\tau$ is the \emph{greatest}
sound assignment.
\end{thm}

\paragraph{Lean.}
\begin{lstlisting}[style=lean]
theorem tau_sound (wf : WellFounded S.Prem) :
    Sound S (tau S wf)
\end{lstlisting}

\paragraph{Haskell.}
\begin{lstlisting}[style=haskell]
-- tau_sound: sound s (tau s)
\end{lstlisting}

%===================================================================
\section{Derivations: max-min semantics agrees with fixpoint}
%===================================================================

\subsection{Inductive type: \lean{Deriv}}

\begin{defn}[Derivation]
A \textbf{derivation} of claim $x$ is a checking step together with
derivations of all its premises (a well-founded tree):
\end{defn}

\paragraph{Lean.}
\begin{lstlisting}[style=lean]
inductive Deriv (S : System C T) : C -> Type where
  | node (f : S.Step) (sub : forall y, y in S.prems f -> Deriv S y) :
      Deriv S (S.concl f)
\end{lstlisting}

\paragraph{Haskell.}
\begin{lstlisting}[style=haskell]
data Deriv c t = Node (Step c t)
                      (Map c (Deriv c t))   -- one sub-derivation per premise
\end{lstlisting}


\subsection{Definition: \lean{foldPrems}}

\begin{defn}[Membership-aware premise fold]
A variant of \lean{cmin}-fold over a list of claims, where the value function
receives a membership proof.  Used to keep \lean{derivGrade} structurally
recursive.
\end{defn}

\paragraph{Lean.}
\begin{lstlisting}[style=lean]
def foldPrems : (l : List C) -> (forall y, y in l -> T) -> T -> T
  | [], _, init => init
  | y :: t, g, init =>
      cmin (g y (.head t))
           (foldPrems t (fun z hz => g z (.tail y hz)) init)
\end{lstlisting}

\paragraph{Haskell.}
\begin{lstlisting}[style=haskell]
foldPrems :: [c] -> (c -> t) -> t -> t
foldPrems [] _ init = init
foldPrems (y:ys) g init = cmin (g y) (foldPrems ys g init)
-- (Membership proof is erased in Haskell.)
\end{lstlisting}


\subsection{Definition: \lean{derivGrade}}

\begin{defn}[Derivation grade]
The \textbf{grade} of a derivation is the meet of every step grade in the
tree --- the weakest link across all checking steps used:
\[
  \mathrm{derivGrade}(\lean{node}\,f\,\mathrm{sub}) =
  \mathrm{foldPrems}\bigl(\mathrm{prems}(f),\;
  y \mapsto \mathrm{derivGrade}(\mathrm{sub}(y)),\;
  \mathrm{grade}(f)\bigr).
\]
\end{defn}

\paragraph{Lean.}
\begin{lstlisting}[style=lean]
noncomputable def derivGrade : {x : C} -> Deriv S x -> T :=
  fun d =>
    Deriv.rec (motive := fun _ _ => T)
      (fun f _ ih =>
        foldPrems (S.prems f) ih (S.grade f)) d
\end{lstlisting}

\paragraph{Haskell.}
\begin{lstlisting}[style=haskell]
derivGrade :: System c t -> Deriv c t -> t
derivGrade s (Node f sub) =
  foldPrems (prems s f) (derivGrade s . sub) (grade s f)
\end{lstlisting}


\subsection{Theorem: \lean{derivGrade\_node}}

\begin{thm}[Computation rule for derivGrade]
\[
  \mathrm{derivGrade}(\lean{node}\,f\,\mathrm{sub}) =
  \mathrm{foldPrems}(\mathrm{prems}(f),\;
  y\,h_y \mapsto \mathrm{derivGrade}(\mathrm{sub}(y, h_y)),\;
  \mathrm{grade}(f)).
\]
\end{thm}

\paragraph{Lean.}
\begin{lstlisting}[style=lean]
theorem derivGrade_node (f : S.Step)
    (sub : forall y, y in S.prems f -> Deriv S y) :
    derivGrade S (.node f sub) =
      foldPrems (S.prems f)
        (fun y hy => derivGrade S (sub y hy)) (S.grade f)
\end{lstlisting}

\paragraph{Haskell.}
\begin{lstlisting}[style=haskell]
-- derivGrade_node: derivGrade (Node f sub) =
--   foldPrems (prems f) (derivGrade . sub) (grade f)
-- This is just the definition; stated as a lemma for rewriting.
\end{lstlisting}


\subsection{Theorems: \lean{foldPrems\_le} and \lean{le\_foldPrems}}

\begin{thm}[foldPrems vs cmin-fold]
These bridge lemmas relate \lean{foldPrems} (which carries membership proofs)
to the ordinary \lean{cmin}-fold over \lean{List.map}:
\begin{align*}
  &\text{If }\forall y\,(h : y \in l),\; g(y,h) \le t(y),\text{ then }
  \mathrm{foldPrems}(l, g, \mathrm{init}) \le (l.\mathrm{map}\,t).\mathrm{foldr}(\mathrm{cmin}, \mathrm{init}).
  && \text{(\lean{foldPrems\_le})}\\
  &\text{If }\forall y\,(h : y \in l),\; t(y) \le g(y,h),\text{ then }
  (l.\mathrm{map}\,t).\mathrm{foldr}(\mathrm{cmin}, \mathrm{init}) \le \mathrm{foldPrems}(l, g, \mathrm{init}).
  && \text{(\lean{le\_foldPrems})}
\end{align*}
\end{thm}

\paragraph{Lean.}
\begin{lstlisting}[style=lean]
theorem foldPrems_le {t : C -> T} {init : T} :
    forall (l : List C) (g : forall y, y in l -> T),
    (forall y (hy : y in l), g y hy <= t y) ->
    foldPrems l g init <= (l.map t).foldr cmin init

theorem le_foldPrems {t : C -> T} {init : T} :
    forall (l : List C) (g : forall y, y in l -> T),
    (forall y (hy : y in l), t y <= g y hy) ->
    (l.map t).foldr cmin init <= foldPrems l g init
\end{lstlisting}

\paragraph{Haskell.}
\begin{lstlisting}[style=haskell]
-- foldPrems_le: g <= t pointwise => foldPrems l g init <= foldr cmin init (map t l)
-- le_foldPrems: t <= g pointwise => foldr cmin init (map t l) <= foldPrems l g init
\end{lstlisting}


\subsection{Theorem: \lean{derivGrade\_le\_tau}}

\begin{thm}[Every derivation's grade is below $\tau$]
For any $x : C$ and derivation $d$ of $x$,
\[
  \mathrm{derivGrade}(d) \le \tau(x).
\]
$\tau$ dominates the max-min semantics from above.
\end{thm}

\paragraph{Lean.}
\begin{lstlisting}[style=lean]
theorem derivGrade_le_tau (wf : WellFounded S.Prem) :
    forall {x : C} (d : Deriv S x),
    derivGrade S d <= tau S wf x
\end{lstlisting}

\paragraph{Haskell.}
\begin{lstlisting}[style=haskell]
-- derivGrade_le_tau: derivGrade d <= tau x
-- for any derivation d of x.
\end{lstlisting}


\subsection{Theorem: \lean{tau\_le\_derivGrade}}

\begin{thm}[$\tau$ is attained by some derivation]
If $\tau(x) \ne \bot$, then there exists a derivation $d$ of $x$ with
$\tau(x) \le \mathrm{derivGrade}(d)$.

Together with \lean{derivGrade\_le\_tau}: $\tau(x)$ equals the supremum of
derivation grades, so the fixpoint form and the derivation-sup form agree.
\end{thm}

\paragraph{Lean.}
\begin{lstlisting}[style=lean]
theorem tau_le_derivGrade (wf : WellFounded S.Prem) :
    forall x : C, tau S wf x != Chain.bot ->
    exists d : Deriv S x, tau S wf x <= derivGrade S d
\end{lstlisting}

\paragraph{Haskell.}
\begin{lstlisting}[style=haskell]
-- tau_le_derivGrade: tau x /= bot =>
--   exists derivation d of x with tau x <= derivGrade d
\end{lstlisting}

%===================================================================
\section{Morphisms and transport of trust}
%===================================================================

\subsection{Structure: \lean{Mor}}

\begin{defn}[Morphism of trust systems]
A \textbf{morphism} $M : S \to S'$ between trust systems (over possibly
different claim types $C, C'$ and chain types $T, T'$) consists of:
\begin{itemize}
  \item $\mathrm{onClaim} : C \to C'$,
  \item $\mathrm{onStep} : S.\lean{Step} \to S'.\lean{Step}$,
  \item $\phi : T \to T'$ (monotone, strict at bottom: $\phi(\bot) = \bot$),
  \item premise compatibility: $S'.\mathrm{prems}(\mathrm{onStep}(f)) =
    \mathrm{prems}(f).\mathrm{map}(\mathrm{onClaim})$,
  \item grade laxity: $\phi(S.\mathrm{grade}(f)) \le
    S'.\mathrm{grade}(\mathrm{onStep}(f))$,
  \item step membership: $f \in \mathrm{stepsInto}(x) \Rightarrow
    \mathrm{onStep}(f) \in \mathrm{stepsInto}'(\mathrm{onClaim}(x))$.
\end{itemize}
The strictness $\phi(\bot) = \bot$ is a correction to the prose
Definition~3.1, discovered by this formalisation: without it, transport
fails on underivable claims.
\end{defn}

\paragraph{Lean.}
\begin{lstlisting}[style=lean]
structure Mor (S : System C T) (S' : System C' T') where
  onClaim  : C -> C'
  onStep   : S.Step -> S'.Step
  phi      : T -> T'
  phi_mono : forall {a b : T}, a <= b -> phi a <= phi b
  phi_bot  : phi Chain.bot = Chain.bot
  prems_eq : forall f, S'.prems (onStep f) = (S.prems f).map onClaim
  grade_le : forall f, phi (S.grade f) <= S'.grade (onStep f)
  steps_mem : forall {x : C} {f : S.Step},
    f in S.stepsInto x -> onStep f in S'.stepsInto (onClaim x)
\end{lstlisting}

\paragraph{Haskell.}
\begin{lstlisting}[style=haskell]
data Mor c t c' t' = Mor
  { onClaim  :: c -> c'
  , onStep   :: Step c t -> Step c' t'
  , phi      :: t -> t'
  -- phi_mono, phi_bot, prems_eq, grade_le, steps_mem are laws
  }
\end{lstlisting}


\subsection{Theorem: \lean{tau\_transport} (Lemma 3.2)}

\begin{thm}[Trust transports laxly along morphisms]
For any morphism $M : S \to S'$ and well-founded premise relations,
\[
  \forall x : C,\quad
  M.\phi(\tau_S(x)) \le \tau_{S'}(M.\mathrm{onClaim}(x)).
\]
Importing along a morphism never claims more than $\phi$ of the source's
canonical trust.
\end{thm}

\paragraph{Lean.}
\begin{lstlisting}[style=lean]
theorem tau_transport {S : System C T} {S' : System C' T'}
    (M : Mor S S')
    (wf : WellFounded S.Prem) (wf' : WellFounded S'.Prem) :
    forall x : C,
    M.phi (tau S wf x) <= tau S' wf' (M.onClaim x)
\end{lstlisting}

\paragraph{Haskell.}
\begin{lstlisting}[style=haskell]
-- tau_transport: phi m (tau s x) <= tau s' (onClaim m x)
-- for any morphism m :: Mor s s'.
\end{lstlisting}

%===================================================================
\section{The circular-citation counterexample (Remark 2.3)}
%===================================================================

\subsection{Inductive type: \lean{Two}}

\begin{defn}[Two-element chain]
The two-level chain $\{\mathtt{lo}, \mathtt{hi}\}$ with $\mathtt{lo} < \mathtt{hi}$,
$\bot = \mathtt{lo}$.
\end{defn}

\paragraph{Lean.}
\begin{lstlisting}[style=lean]
inductive Two where
  | lo
  | hi
deriving DecidableEq

instance : Chain Two where
  le a b := a = .lo \/ b = .hi
  bot := .lo
  -- le_refl, le_trans, le_antisymm, le_total, bot_le, dec all proved
\end{lstlisting}

\paragraph{Haskell.}
\begin{lstlisting}[style=haskell]
data Two = Lo | Hi deriving (Eq, Ord, Show, Bounded)
-- Lo < Hi; bot = Lo.
\end{lstlisting}


\subsection{Definition: \lean{cyc}}

\begin{defn}[Circular citation system]
A trust system on \lean{Bool} claims: the only step into claim $x$ has the
single premise $\neg x$, graded $\mathtt{hi}$.  This encodes a two-claim
cycle: each claim is justified by the other.
\end{defn}

\paragraph{Lean.}
\begin{lstlisting}[style=lean]
def cyc : System Bool Two where
  Step      := Bool
  concl     := id
  prems f   := [!f]
  grade _   := .hi
  stepsInto x := [x]
  mem_stepsInto := ...
  stepsInto_mem f := .head []
\end{lstlisting}

\paragraph{Haskell.}
\begin{lstlisting}[style=haskell]
cyc :: System Bool Two
cyc = System
  { concl = id, prems = \f -> [not f]
  , grade = const Hi, stepsInto = \x -> [x]
  , ... }
\end{lstlisting}


\subsection{Theorem: \lean{cyc\_sound\_top}}

\begin{thm}[Constant-$\mathtt{hi}$ is sound on \lean{cyc}]
The assignment $x \mapsto \mathtt{hi}$ is sound: every node-local certificate
passes.
\end{thm}

\paragraph{Lean.}
\begin{lstlisting}[style=lean]
theorem cyc_sound_top : Sound cyc (fun _ => Two.hi)
\end{lstlisting}

\paragraph{Haskell.}
\begin{lstlisting}[style=haskell]
-- sound cyc (const Hi) = True
-- (Each step is its own certificate: grade hi, premise is the other claim
--  also assigned hi.)
\end{lstlisting}


\subsection{Theorem: \lean{cyc\_no\_deriv}}

\begin{thm}[No claim in \lean{cyc} has any derivation]
\[
  \forall x : \lean{Bool},\; \lean{Deriv}\,\lean{cyc}\,x \to \lean{False}.
\]
The max-min semantics is vacuous: there are no finite derivation trees.
\end{thm}

\paragraph{Lean.}
\begin{lstlisting}[style=lean]
theorem cyc_no_deriv : forall x : Bool, Deriv cyc x -> False
\end{lstlisting}

\paragraph{Haskell.}
\begin{lstlisting}[style=haskell]
-- Any attempt to build a Deriv for cyc loops forever.
\end{lstlisting}


\subsection{Theorem: \lean{cyc\_not\_wf}}

\begin{thm}[\lean{cyc}'s premise relation is not well-founded]
$\neg\,\lean{WellFounded}(\lean{cyc}.\lean{Prem})$.

Therefore the certificate lemma does not apply to \lean{cyc} --- as it must
not, since \lean{cyc\_sound\_top} shows that $\mathtt{hi}$ passes every
local check while \lean{cyc\_no\_deriv} shows no derivation exists.
\end{thm}

\paragraph{Lean.}
\begin{lstlisting}[style=lean]
theorem cyc_not_wf : not WellFounded cyc.Prem
\end{lstlisting}

\paragraph{Haskell.}
\begin{lstlisting}[style=haskell]
-- cyc's dependency graph has a cycle (true -> false -> true -> ...),
-- so the premise relation is not well-founded.
\end{lstlisting}

%===================================================================
\section{Consolidations}
%===================================================================

\subsection{Structure: \lean{Consolidation}}

\begin{defn}[Consolidation]
A \textbf{consolidation} between trust systems $S$ (over $C$) and $S'$ (over
$C'$) over the \emph{same} check-chain $T$ consists of:

\textit{Shapes forward:}
\begin{itemize}
  \item $\mathrm{onClaim} : C \to C'$,
  \item $\mathrm{onStep} : S.\lean{Step} \to S'.\lean{Step}$,
  \item premise compatibility and grade laxity $S.\mathrm{grade}(f) \le
    S'.\mathrm{grade}(\mathrm{onStep}(f))$,
  \item step membership.
\end{itemize}

\textit{Positions backward:} for each $x : C$, every target step
$f' \in S'.\mathrm{stepsInto}(\mathrm{onClaim}(x))$ pulls back to a source
step $\mathrm{posn}(f', h) \in S.\mathrm{stepsInto}(x)$, with premise
correspondence and \emph{citation condition}: $S'.\mathrm{grade}(f') \le
S.\mathrm{grade}(\mathrm{posn}(f', h))$.

The citation condition captures the notion of a cached grade: the consolidated
system records a grade at most as strong as the source's.  Note \lean{posn}
is indexed by the preimage $x$: when consolidation merges claims, every
justification must pull back on each side.
\end{defn}

\paragraph{Lean.}
\begin{lstlisting}[style=lean]
structure Consolidation (S : System C T) (S' : System C' T) where
  -- shapes forward
  onClaim    : C -> C'
  onStep     : S.Step -> S'.Step
  prems_eq   : forall f, S'.prems (onStep f) = (S.prems f).map onClaim
  grade_le   : forall f, S.grade f <= S'.grade (onStep f)
  steps_mem  : forall {x : C} {f : S.Step},
    f in S.stepsInto x -> onStep f in S'.stepsInto (onClaim x)
  -- positions backward
  posn       : forall {x : C} (f' : S'.Step),
    f' in S'.stepsInto (onClaim x) -> S.Step
  posn_mem   : forall {x : C} {f' : S'.Step}
    (h : f' in S'.stepsInto (onClaim x)),
    posn f' h in S.stepsInto x
  posn_prems : forall {x : C} {f' : S'.Step}
    (h : f' in S'.stepsInto (onClaim x)),
    S'.prems f' = (S.prems (posn f' h)).map onClaim
  posn_grade : forall {x : C} {f' : S'.Step}
    (h : f' in S'.stepsInto (onClaim x)),
    S'.grade f' <= S.grade (posn f' h)
\end{lstlisting}

\paragraph{Haskell.}
\begin{lstlisting}[style=haskell]
data Consolidation c t c' = Consolidation
  { onClaim   :: c -> c'
  , onStep    :: Step c t -> Step c' t
  , posn      :: c -> Step c' t -> Step c t
  -- prems_eq, grade_le, steps_mem, posn_mem, posn_prems, posn_grade are laws
  }
\end{lstlisting}

\begin{fidelitynote}
The \lean{posn} field takes a membership proof $h : f' \in
\mathrm{stepsInto}'(\mathrm{onClaim}(x))$ as an argument.  This is the
dependent type that records \emph{which} preimage $x$ we are pulling back to.
When consolidation merges multiple claims, \lean{posn} must be defined
separately for each preimage; the Haskell approximation adds $x$ as an extra
argument to \lean{posn}, but loses the membership-proof dependency.
\end{fidelitynote}


\subsection{Definition: \lean{Consolidation.toMor}}

\begin{defn}[Forward half of consolidation is a morphism]
The forward data of a consolidation $K : S \to S'$ at $\phi = \mathrm{id}$
assembles into a morphism \lean{Mor S S'}.
\end{defn}

\paragraph{Lean.}
\begin{lstlisting}[style=lean]
def Consolidation.toMor {S : System C T} {S' : System C' T}
    (K : Consolidation S S') : Mor S S' where
  onClaim   := K.onClaim
  onStep    := K.onStep
  phi       := id
  phi_mono  := fun h => h
  phi_bot   := rfl
  prems_eq  := K.prems_eq
  grade_le  := K.grade_le
  steps_mem := K.steps_mem
\end{lstlisting}

\paragraph{Haskell.}
\begin{lstlisting}[style=haskell]
toMor :: Consolidation c t c' -> Mor c t c' t
toMor k = Mor (onClaim k) (onStep k) id
\end{lstlisting}


\subsection{Theorem: \lean{Consolidation.tau\_le}}

\begin{thm}[Position map forbids trust inflation]
A consolidation $K : S \to S'$ satisfies
\[
  \forall x : C,\quad \tau_{S'}(\mathrm{onClaim}(x)) \le \tau_S(x).
\]
The position map pulls every target justification back to a source
justification, preventing the consolidated system from claiming more than
the source.
\end{thm}

\paragraph{Lean.}
\begin{lstlisting}[style=lean]
theorem Consolidation.tau_le {S : System C T} {S' : System C' T}
    (K : Consolidation S S')
    (wf : WellFounded S.Prem) (wf' : WellFounded S'.Prem) :
    forall x : C, tau S' wf' (K.onClaim x) <= tau S wf x
\end{lstlisting}

\paragraph{Haskell.}
\begin{lstlisting}[style=haskell]
-- tau_le: tau s' (onClaim k x) <= tau s x
-- for all x, given consolidation k :: Consolidation s s'.
\end{lstlisting}


\subsection{Theorem: \lean{tau\_consolidate} (Theorem 3)}

\begin{thm}[Consolidation preservation (Theorem 3)]
A consolidation $K : S \to S'$ with a genuine position map preserves
canonical trust on the nose:
\[
  \forall x : C,\quad \tau_{S'}(\mathrm{onClaim}(x)) = \tau_S(x).
\]
Forward inequality: shapes transport trust laxly (\lean{tau\_transport} at
$\phi = \mathrm{id}$).  Backward inequality: positions pull every target
justification back (\lean{tau\_le}).
\end{thm}

\paragraph{Lean.}
\begin{lstlisting}[style=lean]
theorem tau_consolidate {S : System C T} {S' : System C' T}
    (K : Consolidation S S')
    (wf : WellFounded S.Prem) (wf' : WellFounded S'.Prem) :
    forall x : C,
    tau S' wf' (K.onClaim x) = tau S wf x
\end{lstlisting}

\paragraph{Haskell.}
\begin{lstlisting}[style=haskell]
-- tau_consolidate: tau s' (onClaim k x) == tau s x
-- for all x, given a full consolidation k (with position map).
\end{lstlisting}

%===================================================================
\section{Composition of consolidations}
%===================================================================

\subsection{Definition: \lean{Consolidation.id}}

\begin{defn}[Identity consolidation]
The identity consolidation on $S$: shapes and positions both identity, grades
exact ($S.\mathrm{grade}(f) \le S.\mathrm{grade}(f)$ by reflexivity).
\end{defn}

\paragraph{Lean.}
\begin{lstlisting}[style=lean]
protected def Consolidation.id (S : System C T) :
    Consolidation S S where
  onClaim := id
  onStep  := id
  prems_eq f := (List.map_id (S.prems f)).symm
  grade_le f := Chain.le_refl (S.grade f)
  steps_mem h := h
  posn f' _ := f'
  posn_mem h := h
  posn_prems {_ f'} _ := (List.map_id (S.prems f')).symm
  posn_grade {_ f'} _ := Chain.le_refl (S.grade f')
\end{lstlisting}

\paragraph{Haskell.}
\begin{lstlisting}[style=haskell]
consolidationId :: Consolidation c t c
consolidationId = Consolidation id id (\_ f' -> f')
\end{lstlisting}


\subsection{Definition: \lean{Consolidation.comp}}

\begin{defn}[Composition of consolidations]
Given $K : S \to S'$ and $L : S' \to S''$, their composite $K \cdot L :
S \to S''$ has:
\begin{itemize}
  \item shapes: $\mathrm{onClaim} = L.\mathrm{onClaim} \circ K.\mathrm{onClaim}$,
    $\mathrm{onStep} = L.\mathrm{onStep} \circ K.\mathrm{onStep}$;
  \item positions: $\mathrm{posn}(f'', h) =
    K.\mathrm{posn}(L.\mathrm{posn}(f'', h),\, L.\mathrm{posn\_mem}(h))$
    (contravariant --- pull back through $L$ first, then through $K$);
  \item grades: by transitivity of $\le$.
\end{itemize}
\end{defn}

\paragraph{Lean.}
\begin{lstlisting}[style=lean]
protected def Consolidation.comp
    {S : System C T} {S' : System C' T} {S'' : System C'' T}
    (K : Consolidation S S') (L : Consolidation S' S'') :
    Consolidation S S'' where
  onClaim x := L.onClaim (K.onClaim x)
  onStep f  := L.onStep (K.onStep f)
  prems_eq f := by rw [L.prems_eq, K.prems_eq, List.map_map]; rfl
  grade_le f := Chain.le_trans (K.grade_le f) (L.grade_le (K.onStep f))
  steps_mem h := L.steps_mem (K.steps_mem h)
  posn f'' h  := K.posn (L.posn f'' h) (L.posn_mem h)
  posn_mem h  := K.posn_mem (L.posn_mem h)
  posn_prems h := by
    rw [L.posn_prems h, K.posn_prems (L.posn_mem h), List.map_map]; rfl
  posn_grade h := Chain.le_trans (L.posn_grade h)
                    (K.posn_grade (L.posn_mem h))
\end{lstlisting}

\paragraph{Haskell.}
\begin{lstlisting}[style=haskell]
consolidationComp :: Consolidation c t c'
                  -> Consolidation c' t c''
                  -> Consolidation c t c''
consolidationComp k l = Consolidation
  { onClaim = onClaim l . onClaim k
  , onStep  = onStep l . onStep k
  , posn    = \x f'' -> posn k (onClaim k x) (posn l (onClaim l (onClaim k x)) f'')
  }
\end{lstlisting}


\subsection{Theorem: \lean{tau\_consolidate\_comp}}

\begin{thm}[End-to-end preservation through a pipeline]
For composable consolidations $K : S \to S'$ and $L : S' \to S''$:
\[
  \forall x : C,\quad
  \tau_{S''}((K \cdot L).\mathrm{onClaim}(x)) = \tau_S(x).
\]
This factors as $\tau_{S''}(L.\mathrm{onClaim}(K.\mathrm{onClaim}(x)))
= \tau_{S'}(K.\mathrm{onClaim}(x)) = \tau_S(x)$, so end-to-end trust
preservation is the conjunction of the per-layer equations.
\end{thm}

\paragraph{Lean.}
\begin{lstlisting}[style=lean]
theorem tau_consolidate_comp
    {S : System C T} {S' : System C' T} {S'' : System C'' T}
    (K : Consolidation S S') (L : Consolidation S' S'')
    (wf : WellFounded S.Prem)
    (wf' : WellFounded S'.Prem)
    (wf'' : WellFounded S''.Prem) :
    forall x : C,
    tau S'' wf'' ((K.comp L).onClaim x) = tau S wf x
\end{lstlisting}

\paragraph{Haskell.}
\begin{lstlisting}[style=haskell]
-- tau_consolidate_comp:
--   tau s'' (onClaim (comp k l) x) == tau s x
-- for all x, composable consolidations k, l.
\end{lstlisting}

%===================================================================
\section{The shape-only counterexample}
%===================================================================

\subsection{Definitions: \lean{ground} and \lean{layer4}}

\begin{defn}[Shape-only counterexample]
Two systems over \lean{Unit} claims, with the two-element chain \lean{Two}:

\textbf{\lean{ground}}: one claim, one $\mathtt{hi}$-graded ground check
(no premises).
\[
  \lean{Step} = \lean{Unit},\;
  \mathrm{prems} = [],\;
  \mathrm{grade} = \mathtt{hi},\;
  \mathrm{stepsInto}(\langle\rangle) = [\langle\rangle].
\]

\textbf{\lean{layer4}}: same claim, no steps at all (the shape-only
consolidated version: the claim survives but the check is dropped).
\[
  \lean{Step} = \lean{Empty},\;
  \mathrm{stepsInto}(\langle\rangle) = [].
\]
\end{defn}

\paragraph{Lean.}
\begin{lstlisting}[style=lean]
def ground : System Unit Two where
  Step := Unit
  concl _ := ()
  prems _ := []
  grade _ := .hi
  stepsInto _ := [()]
  mem_stepsInto _ := rfl
  stepsInto_mem _ := .head []

def layer4 : System Unit Two where
  Step := Empty
  concl := Empty.elim
  prems := Empty.elim
  grade := Empty.elim
  stepsInto _ := []
  mem_stepsInto := fun h => nomatch h
  stepsInto_mem := fun f => f.elim
\end{lstlisting}

\paragraph{Haskell.}
\begin{lstlisting}[style=haskell]
-- ground: one ground check graded Hi; layer4: no steps.
ground :: System () Two
ground = System { step=[()] , concl=const ()
                , prems=const [], grade=const Hi
                , stepsInto=const [()] }

data Void  -- empty type
layer4 :: System () Two
layer4 = System { step=[], concl=absurd
                , prems=absurd, grade=absurd
                , stepsInto=const [] }
\end{lstlisting}


\subsection{Theorems: \lean{ground\_wf} and \lean{layer4\_wf}}

\begin{thm}[Well-foundedness of both systems]
Both \lean{ground} and \lean{layer4} have well-founded premise relations
(both empty).
\end{thm}

\paragraph{Lean.}
\begin{lstlisting}[style=lean]
theorem ground_wf : WellFounded ground.Prem
theorem layer4_wf : WellFounded layer4.Prem
\end{lstlisting}

\paragraph{Haskell.}
\begin{lstlisting}[style=haskell]
-- Both systems have empty premise relations; trivially well-founded.
\end{lstlisting}


\subsection{Theorems: \lean{ground\_tau\_hi} and \lean{layer4\_tau\_bot}}

\begin{thm}[Trust degrades from \lean{hi} to \lean{bot} after shape-only consolidation]
\[
  \tau_{\lean{ground}}(\langle\rangle) = \mathtt{hi},
  \qquad
  \tau_{\lean{layer4}}(\langle\rangle) = \mathtt{lo} = \bot.
\]
This is the formal shadow of the \texttt{synthesis-layer4.md} failure: after
four shape-only consolidation layers, machine-resolvable citations are gone,
and canonical trust collapses to $\bot$.
\end{thm}

\paragraph{Lean.}
\begin{lstlisting}[style=lean]
theorem ground_tau_hi :
    tau ground ground_wf () = Two.hi

theorem layer4_tau_bot :
    tau layer4 layer4_wf () = Two.lo
\end{lstlisting}

\paragraph{Haskell.}
\begin{lstlisting}[style=haskell]
-- tau ground () == Hi    (ground check graded Hi, no premises)
-- tau layer4 () == Lo    (no steps -> foldr cmax bot [] = bot)
\end{lstlisting}

%===================================================================
\section{Translation conventions}
%===================================================================

\begin{enumerate}
\item \textbf{Dependent types throughout.}  \lean{System.Step} varies between
  instances; \lean{Consolidation.posn} depends on a membership proof;
  \lean{Deriv S x} is indexed by the claim $x$.  All of these are
  approximated by erasing the dependency in Haskell.

\item \textbf{Noncomputable definitions.}  \lean{tau} and \lean{derivGrade}
  are marked \lean{noncomputable} because they use \lean{Classical.choice}
  (via \lean{WellFounded.fix} and \lean{Deriv.rec}).  In Haskell they are
  written as recursive functions; termination is assumed.

\item \textbf{The three roles of linearity.}  The \lean{Chain} hypothesis
  (totality) is used at exactly three points: (a) \lean{cmax\_cases}
  (attainment in \lean{tau\_sound}), (b) \lean{foldr\_max\_cases}
  (attainment in \lean{tau\_le\_derivGrade}), and (c) antisymmetry in
  \lean{tau\_le\_derivGrade}'s uniqueness argument.  All three fail for
  general lattices.

\item \textbf{Phi-bot as a correction.}  The \lean{phi\_bot} field of
  \lean{Mor} ($\phi(\bot) = \bot$) is not in the original prose definition.
  Without it, \lean{tau\_transport} fails for claims with no derivation
  (where $\tau = \bot$) if $\phi(\bot) > \bot$ in the target chain.  This
  is a genuine mathematical correction discovered by Lean.

\item \textbf{Consolidation vs morphism.}  A \lean{Consolidation} is a
  \lean{Mor} at $\phi = \mathrm{id}$ plus a backward position map.
  The extra structure (positions backward) is what yields preservation
  ($=$) rather than lax transport ($\le$).

\item \textbf{Proof bodies omitted.}  Lean checks all proofs.
\end{enumerate}

\end{document}
